%% ============================================================
%%  SOLUCIÓN — Ejercicio 4: Artículo básico completo (INTEGRADOR)
%%  Clase 1 — Al final de la Clase 1
%%  Compilar: F5 (pdflatex, dos veces para TOC y referencias)
%% ============================================================

\documentclass[12pt, a4paper]{article}

% --- Codificación y lenguaje ---
\usepackage[utf8]{inputenc}
\usepackage[T1]{fontenc}
\usepackage[spanish]{babel}
\usepackage{lmodern}

% --- Márgenes ---
\usepackage[top=2.5cm, bottom=2.5cm, left=3cm, right=3cm]{geometry}

% --- Matemáticas y tablas ---
\usepackage{amsmath, amssymb}
\usepackage{booktabs}
\usepackage{graphicx}
\usepackage[colorlinks=true, linkcolor=black]{hyperref}

% --- Metadatos ---
\title{Salarios y educación en Paraguay}
\author{Tu Nombre}
\date{\today}

% ================================================================
\begin{document}
% ================================================================

\maketitle
\thispagestyle{empty}

% --- Abstract ---
\begin{abstract}
\noindent
Este trabajo analiza el efecto causal de la educación sobre los salarios 
de los trabajadores paraguayos, utilizando datos de la Encuesta Permanente 
de Hogares 2023. Encontramos que un año adicional de educación está asociado 
con un aumento del 6.9\% en el salario mensual. Las mujeres, a pesar de tener 
niveles educativos similares, ganan en promedio 18.3\% menos que los hombres. 
Estos resultados sugieren que la inversión en educación es rentable, pero que 
existen importantes brechas de género en el mercado laboral que requieren 
atención de política pública.
\end{abstract}

\newpage
\tableofcontents
\newpage

% ================================================================
\section{Introducción}
% ================================================================

Entender los determinantes de los salarios es una pregunta central en economía 
laboral. Desde el trabajo seminal de Mincer (1974), sabemos que los salarios 
crecen con la educación y la experiencia. Sin embargo, poco se conoce sobre 
estos retornos en el contexto paraguayo.

Este trabajo estima ecuaciones de ingresos para Paraguay usando datos de la 
Encuesta Permanente de Hogares 2023. Nuestro objetivo es cuantificar el 
retorno a la educación y documentar disparidades de género en los salarios.

\subsection{Motivación}

Comprender los retornos a la educación es \emph{fundamental} para el diseño 
de políticas públicas educativas. Si los retornos son altos, entonces invertir 
en educación es rentable no solo para el individuo, sino para la sociedad en 
su conjunto. Si existen disparidades de género, eso sugiere que hay factores 
no educativos (discriminación, segregación ocupacional) afectando los salarios.

% ================================================================
\section{El modelo econométrico}
% ================================================================

Nuestro modelo base es la ecuación de Mincer modificada:
\begin{equation}
  \ln(w_i) = \alpha + \beta_1 \, \text{educ}_i + \beta_2 \, \text{exp}_i 
             + \beta_3 \, \text{exp}_i^2 + \varepsilon_i
  \label{eq:mincer}
\end{equation}

donde $w_i$ es el salario mensual del individuo $i$, $\text{educ}_i$ son los 
años de escolaridad, $\text{exp}_i$ es la experiencia potencial (edad $-$ años 
de educación $-$ 6), y $\varepsilon_i$ es el término de error.

El parámetro $\beta_1$ en la ecuación \eqref{eq:mincer} se interpreta como el 
retorno porcentual a un año adicional de educación. Los parámetros $\beta_2$ 
y $\beta_3$ capturan el ciclo de vida de los ingresos, que típicamente crece 
con la experiencia pero a tasa decreciente.

% ================================================================
\section{Datos}
% ================================================================

Los datos provienen de la Encuesta Permanente de Hogares 2023, realizada por 
el Instituto Nacional de Estadística. La muestra cubre 1.200 individuos 
ocupados entre 18 y 65 años.

\begin{table}[htbp]
  \centering
  \caption{Estadísticas descriptivas de la muestra}
  \label{tab:desc}
  \begin{tabular}{lrrrr}
    \toprule
    Variable                & Media  & D.E.   & Mín. & Máx. \\
    \midrule
    Salario (Gs.)           & 2.450.000 & 1.120.000 & 550.000 & 9.800.000 \\
    Log(salario)            & 14.52  & 0.53   & 13.22 & 16.10 \\
    Años de educación       & 12.3   & 3.4    & 0    & 20   \\
    Experiencia (años)      & 8.7    & 5.1    & 0    & 40   \\
    Mujer (indicador)       & 0.48   & 0.50   & 0    & 1    \\
    \midrule
    Observaciones           & \multicolumn{4}{c}{1.200} \\
    \bottomrule
  \end{tabular}
  \begin{minipage}{\linewidth}
    \footnotesize \textit{Notas:} El salario en guaraníes fue deflactado 
    al año base 2018. La experiencia se calcula como edad $-$ años de 
    educación $-$ 6. \textit{Fuente:} EPH 2023, INE.
  \end{minipage}
\end{table}

La Tabla \ref{tab:desc} muestra que el salario promedio es de aproximadamente 
2.45 millones de guaraníes mensuales, con considerable variabilidad (D.E. = 
1.12M). El nivel educativo promedio es de 12.3 años (poco más que primaria 
completa + secundaria incompleta).

% ================================================================
\section{Resultados}
% ================================================================

En esta sección presentamos los resultados de estimaciones OLS de la ecuación 
de Mincer bajo distintas especificaciones.

\subsection{Modelo base}

La primera columna de la Tabla \ref{tab:reg} es el modelo base sin controles 
adicionales. El coeficiente de educación es 0.082***, lo que significa que 
cada año adicional de educación está asociado con un aumento de 8.2\% en el 
salario.

\subsection{Con controles de experiencia}

Cuando incluimos la experiencia y su cuadrado (columna 2), el efecto de la 
educación disminuye ligeramente a 0.071***. Esto sugiere que educación y 
experiencia están parcialmente correlacionadas, pero los retornos a la 
educación siguen siendo importantes.

\subsection{Con control de género}

La especificación 3 incluye un indicador de género (mujer = 1). Encontramos 
una brecha de género significativa: las mujeres ganan 18.3\% menos que los 
hombres, incluso después de controlar por educación y experiencia.

\begin{table}[htbp]
  \centering
  \caption{Estimaciones OLS: Ecuación de Mincer}
  \label{tab:reg}
  \begin{tabular}{lccc}
    \toprule
                    & (1)        & (2)        & (3)        \\
                    & \multicolumn{3}{c}{\textit{Var. dep.: } $\ln(\text{salario})$} \\
    \midrule
    Educación       & 0.082***   & 0.071***   & 0.069***   \\
                    & (0.008)    & (0.009)    & (0.010)    \\[4pt]
    Experiencia     &            & 0.034**    & 0.031**    \\
                    &            & (0.012)    & (0.013)    \\[4pt]
    Experiencia$^2$ &            & $-$0.001** & $-$0.001*  \\
                    &            & (0.000)    & (0.000)    \\[4pt]
    Mujer (=1)      &            &            & $-$0.183***\\
                    &            &            & (0.042)    \\[4pt]
    Constante       & 7.241***   & 7.018***   & 7.312***   \\
                    & (0.098)    & (0.115)    & (0.122)    \\
    \midrule
    Observaciones   & 1.200      & 1.200      & 1.200      \\
    $R^2$ ajustado  & 0.241      & 0.318      & 0.374      \\
    \bottomrule
  \end{tabular}
  \begin{minipage}{\linewidth}
    \footnotesize \textit{Notas:} Errores estándar robustos a 
    heterocedasticidad en paréntesis. *** $p<0.01$, ** $p<0.05$, 
    * $p<0.1$. \textit{Fuente:} EPH 2023, elaboración propia.
  \end{minipage}
\end{table}

% ================================================================
\section{Conclusiones}
% ================================================================

Este trabajo documenta que los retornos a la educación en Paraguay son 
sustanciales, del orden de 7\% por año de escolaridad adicional. Sin embargo, 
también encontramos una importante brecha de género que persiste incluso después 
de controlar por educación y experiencia.

Estos resultados sugieren que:
\begin{enumerate}
  \item La inversión en educación es rentable para los individuos.
  \item Existen factores más allá de la educación que explican los salarios 
        (discriminación, capital social, acceso a empleos de calidad).
  \item Las políticas públicas deben no solo expandir la cobertura educativa, 
        sino también abordar las barreras que enfrentan las mujeres en el 
        mercado laboral.
\end{enumerate}

% ================================================================
\end{document}
